
\section{Inverse-IFEval Results}
\label{sec:appendix:lab}

\subsection{Results by cover size}
\label{sec:appendix:strata}

\begin{table*}[t]
  \centering\small
  \begin{tabularx}{\textwidth}{@{}>{\raggedright\arraybackslash}Xrrr@{}}
    \toprule
    Arm & $N$ & Excess size (\%) & Sat.\ (\%) \\
    \midrule
    \multicolumn{4}{@{}l}{\emph{$|D_\star| = 2$}} \\
    no prompt comments & 423 & $+74.2$ $[+52.2, +98.8]$ & 38.7 \\
    comment-shaped noise & 423 & $+65.1$ $[+48.2, +84.4]$ & 39.5 \\
    informative comments & 423 & $-1.1$ $[-9.2, +7.7]$ & 39.0 \\
    \midrule
    \multicolumn{4}{@{}l}{\emph{$|D_\star| = 3$}} \\
    no prompt comments & 129 & $+15.2$ $[+0.5, +31.8]$ & 40.3 \\
    comment-shaped noise & 129 & $+14.2$ $[-1.8, +32.6]$ & 42.9 \\
    informative comments & 129 & $-21.2$ $[-28.7, -13.2]$ & 35.7 \\
    \midrule
    \multicolumn{4}{@{}l}{\emph{pooled}} \\
    no prompt comments & 552 & $+60.4$ $[+43.8, +80.7]$ & 39.0 \\
    comment-shaped noise & 552 & $+53.2$ $[+39.4, +68.7]$ & 40.3 \\
    informative comments & 552 & $-5.8$ $[-12.0, +1.1]$ & 38.2 \\
    \bottomrule
  \end{tabularx}
  \caption{
    The ratchet holds in both strata, and the protocol reverses it in both. Each row is an arm
    within one $|D_\star|$ stratum, with pooled rows below. The columns give units, final excess
    size $\frac{|D_T|}{|D_\star|}{-}1$ with its 95\% per-world percentile-bootstrap CI, and
    last-3-step constraint satisfaction. The uncommented arm's CI excludes cover in both strata,
    far above it where the cover is two instructions and by a smaller margin where it is three, so
    the pooled result is not carried by one slice. In the three-instruction stratum the protocol
    arm prunes \emph{below} cover, a cost \S\ref{sec:appendix:calibration} works out. Covers here
    run to at most three instructions, an order of magnitude below the median context file of
    \S\ref{sec:growth}.
  }
  \label{tab:strata}
\end{table*}


\autoref{tab:strata} is the full split behind \autoref{sec:protocol}'s pooled numbers, and the two
strata disagree about the headline. At covers of two, the uncommented arm ends far above cover. At
covers of three, its interval straddles cover, so criterion~(i) of \autoref{tab:seeds} does not
hold within the three-constraint slice on its own. The pooled ratchet is therefore a two-constraint
result at this horizon.

The disagreement shows \autoref{sec:theory}'s interference regime appearing inside the experiment
rather than a defect in it. Interference makes deletion pay, and it is weak at these constraint
counts and strong at the counts \autoref{sec:growth} measures. Our scope is therefore split: the
experiment demonstrates the mechanism where retention is nearly free, and the corpus shows the
ratchet running where pruning would pay.

\subsection{The ratchet against maintainer capability}
\label{sec:appendix:capacity}

\begin{table*}[t]
  \centering\small
  \begin{tabularx}{\textwidth}{@{}>{\raggedright\arraybackslash}Xrrrrrr@{}}
    \toprule
    & \multicolumn{3}{c}{Excess size} & \multicolumn{3}{c}{Satisfaction} \\
    \cmidrule(lr){2-4} \cmidrule(lr){5-7}
    Maintainer & Control & $\Delta$ \textsc{commit} (\%) & $z$ & Control (\%) & $\Delta$ \textsc{commit} (\%) & $z$ \\
    \midrule
    Haiku 4.5 & $1.68$ & $-42.1$ & $-2.45$ & $40.9$ & $-4.5$ & $-0.72$ \\
    Sonnet 5 & $3.07$ & $\mathbf{-61.7}$ & $-8.40$ & $38.5$ & $+7.4$ & $0.89$ \\
    Opus 5 & $6.72$ & $-52.3$ & $-9.82$ & $45.3$ & $\mathbf{+18.4}$ & $2.12$ \\
    \bottomrule
  \end{tabularx}
  \caption{
    Stronger maintainers ratchet harder, and the protocol helps them more. Rows
    share one executor (Haiku 4.5), task stream, horizon and protocol; only the maintainer moves.
    Each block gives the uncommented arm's own level --- final size $|D_T|/|D_\star|$,
    last-3-step constraint satisfaction --- then \textsc{commit} against that same row's control,
    since a maintainer swap moves control too. $z$ is the paired per-world statistic;
    \textbf{bold} is the largest $\Delta$ in its block among rows where $|z|$ resolves it.
    Sonnet~5 overran the 1{,}024-character comment budget on two of every five additions.
    $n{=}64$ worlds, one seed.
  }
  \label{tab:capacity}
\end{table*}


\autoref{tab:capacity} moves the maintainer alone, holding the executor, task stream, horizon and
protocol fixed, and answers whether the effect needs a matched maintainer/executor pair. It does not:
the uncommented arm's excess grows with capability, and the comment protocol's advantage grows with it.
The Sonnet~5 row is reported for completeness but is a compliance record rather than a capability
reading --- that maintainer overran the comment budget on two of every five additions, so it did not
write the protocol it was given.

\subsection{What the comment channel needs}
\label{sec:appendix:ablations}

\begin{table*}[t]
  \centering\small
  \begin{tabularx}{\textwidth}{@{}>{\raggedright\arraybackslash}Xrrrr@{}}
    \toprule
    & \multicolumn{2}{c}{Excess size} & \multicolumn{2}{c}{Satisfaction} \\
    \cmidrule(lr){2-3} \cmidrule(lr){4-5}
    Ablation & $\Delta$ (\%) & $z$ & $\Delta$ (\%) & $z$ \\
    \midrule
    comment-shaped noise & $-3.1$ & $-0.33$ & $+3.3$ & $0.77$ \\
    narrative without outcomes & $+8.8$ & $0.86$ & $-3.3$ & $-0.82$ \\
    \textsc{commit}, full schema & $-48.9$ & $-2.78$ & $+4.1$ & $0.50$ \\
    \quad -- round counter & $-47.4$ & $-2.63$ & $-1.9$ & $-0.23$ \\
    \quad -- recurrence count & $-30.9$ & $-1.83$ & $-3.2$ & $-0.42$ \\
    \quad -- falsification lineage & $-34.3$ & $-2.14$ & $-9.6$ & $-1.19$ \\
    \quad -- verbatim restore quote & $-43.5$ & $-2.43$ & $-8.9$ & $-1.09$ \\
    \bottomrule
  \end{tabularx}
  \caption{
    A comment must carry outcomes. Every row is that run's arm against its own control.
    Comment-shaped noise carries control's information at the comment arm's surface form; the
    clause rows drop each of the \textsc{commit} schema's four fields in turn. $z$ is the paired
    per-world statistic on the $\Delta$ beside it. $n{=}64$ worlds and one seed for the clause
    rows; the two content rows are full-pool runs.
  }
  \label{tab:ablations}
\end{table*}


\autoref{tab:ablations} separates three objections: whether the effect is the comment's text or its
content, whether the rationale needs a channel of its own, and which field of the schema carries the
pruning. Rows are each run's arm against its own uncommented control. The clause rows drop one
field of the \textsc{commit} schema each, under the names \autoref{fig:renderings}(a) gives them;
each variant is otherwise byte-identical to the full protocol, so a row's contrast is the field it
drops and not a rewording.

\subsection{Replication across world draws}
\label{sec:appendix:seeds}

\begin{table*}[t]
  \centering\small
  \begin{tabularx}{\textwidth}{@{}>{\raggedright\arraybackslash}Xrrrrr@{}}
    \toprule
    Excess size (\%) & s0 & s1 & s2 & Pooled & sd \\
    \midrule
    no prompt comments & $+59.9$ & $+49.8$ & $+71.6$ & $+60.4$ & 10.9 \\
    comment-shaped noise & $+54.9$ & $+46.0$ & $+58.8$ & $+53.2$ & 6.5 \\
    informative comments & $-0.1$ & $-2.1$ & $-15.1$ & $-5.8$ & 8.2 \\
    \midrule
    Satisfaction (\%) & s0 & s1 & s2 & Pooled & sd \\
    \midrule
    no prompt comments & $38.8$ & $38.3$ & $40.0$ & $39.0$ & 0.9 \\
    comment-shaped noise & $40.0$ & $40.0$ & $40.8$ & $40.3$ & 0.4 \\
    informative comments & $38.9$ & $36.8$ & $39.0$ & $38.2$ & 1.3 \\
    \midrule
    Criterion of record & s0 & s1 & s2 & Pooled & \\
    \midrule
    (i) uncommented arm ratchets & $\checkmark$ & $\checkmark$ & $\checkmark$ & $\checkmark$ & \\
    (iv) noise $\approx$ uncommented & $\checkmark$ & $\checkmark$ & $\checkmark$ & $\checkmark$ & \\
    \bottomrule
  \end{tabularx}
  \caption{
    Every criterion passes on every world draw. The upper block gives final excess size
    $\frac{|D_T|}{|D_\star|}{-}1$ per arm, with one column per seed, then the pooled value and the
    across-seed standard deviation in pp. Each seed is a fresh permutation sample of IFEval
    items, so the seeds are independent draws of the world pool and not repeated runs on one pool.
    The middle block gives last-three-round constraint satisfaction as a level, so the size result
    is readable beside the correctness it was bought at. The lower block gives the criteria of
    record. We fixed these before the run and evaluate
    each one per seed, since a criterion that passed pooled but failed on one seed of three would
    be an averaging artifact, and a gap here would show it. Across seeds, the uncommented arm varies little
    against the effect this design is powered on, which is why we read the pooled numbers in
    \S\ref{sec:protocol}.
  }
  \label{tab:seeds}
\end{table*}


All four criteria of record hold on each of the three world draws separately (\autoref{tab:seeds}),
not only pooled. Each seed is a fresh permutation sample of the eligible IFEval pool, so the three
columns are independent draws of the world pool rather than repeated runs on one. The uncommented
arm's across-seed spread of \LabControlSeedSdPp{} is small against the \LabPracticalDeltaExcessPp{}
effect the design is powered on.

\subsection{What the arm mean hides}
\label{sec:appendix:calibration}

\begin{table*}[t]
  \centering\small
  \begin{tabular}{@{}lrrrrrr@{}}
    \toprule
    \multicolumn{7}{@{}l}{\textbf{(a)} Per-world calibration to the minimum cover} \\
    \midrule
    Arm & $N$ & Mean ratio & Mean $|{\rm ratio}-1|$ & Empty & Within $\frac{1}{4}$ & Above $1.5\times$ \\
    no prompt comments & 552 & 1.604 & 0.921 & 2.4\% & 33.3\% & 22.3\% \\
    comment-shaped noise & 552 & 1.532 & 0.822 & 1.1\% & 35.0\% & 21.9\% \\
    informative comments & 552 & 0.942 & 0.481 & 12.1\% & 37.3\% & 9.8\% \\
    \midrule
    \multicolumn{7}{@{}l}{\textbf{(b)} Satisfaction on the worlds the protocol emptied, and on the rest} \\
    \midrule
    Arm & $N$ & Satisfaction & 95\% CI & \multicolumn{2}{l}{Paired $\Delta$ vs.\ the protocol arm} & 95\% CI \\
    \multicolumn{7}{@{}l}{\emph{worlds the protocol emptied}} \\
    no prompt comments & 67 & 0.279 & $[0.209, 0.348]$ & \multicolumn{2}{l}{$-0.020$} & $[-0.065, +0.025]$ \\
    comment-shaped noise & 67 & 0.328 & $[0.256, 0.400]$ & \multicolumn{2}{l}{$-0.070$} & $[-0.114, -0.027]$ \\
    informative comments & 67 & 0.259 & $[0.199, 0.321]$ & \multicolumn{2}{l}{---} &  \\
    \multicolumn{7}{@{}l}{\emph{worlds it kept instructions on}} \\
    no prompt comments & 485 & 0.406 & $[0.377, 0.436]$ & \multicolumn{2}{l}{$-0.007$} & $[-0.031, +0.016]$ \\
    comment-shaped noise & 485 & 0.413 & $[0.383, 0.443]$ & \multicolumn{2}{l}{$-0.014$} & $[-0.034, +0.007]$ \\
    informative comments & 485 & 0.399 & $[0.368, 0.429]$ & \multicolumn{2}{l}{---} &  \\
    \bottomrule
  \end{tabular}
  \caption{
    The commented arm is closest to cover on the mean and on individual prompts alike.
    \textbf{(a)} gives per-world calibration at the final step: mean $|D_T|/|D_\star|$, mean
    $\bigl||D_T|/|D_\star| - 1\bigr|$ (0 = every world exactly at cover), and the shares of worlds
    ending empty, within a quarter of cover, and above $1.5\times$ it. \emph{informative comments}
    wins on both readings, so its arm mean can be read directly: the two need not agree, and where
    they disagree the mean is the one that misleads. The mean still hides a tail. One prompt in
    eight ends empty under the protocol against one in forty under the uncommented arm.
    \textbf{(b)} prices that tail. We split worlds by whether the protocol arm ended empty, then
    give last-3-step satisfaction per arm within each split and the paired difference from the
    protocol arm on the same worlds, with its 95\% bootstrap CI. Parity holds where the protocol
    kept instructions and fails where it emptied them. Those are the hard worlds, and every arm
    scores far below its own average on them. The protocol's own endpoint defines the split, so
    this paired deficit bounds the cost of deletion rather than estimating it.
  }
  \label{tab:calibration}
\end{table*}


The protocol reaches cover per world and not only on the arm mean
(\autoref{tab:calibration}a), and it overshoots on the short side: it empties one prompt in eight,
against almost none under either uninformed arm. \autoref{tab:calibration}(b) prices that. We split
worlds by whether the protocol ended empty and read every arm within each split, which pairs the
comparison on the same worlds and holds difficulty fixed. Satisfaction parity holds where the
protocol kept instructions and fails where it emptied them. Those are the hard worlds, where every
arm scores far below its average, and the protocol's own endpoint defines the split, so the paired
deficit bounds the cost of deletion rather than estimating it.

One condition on the recommendation follows. Writing a comment risks nothing, because it removes
nothing. Acting on one to delete is where the risk sits, and this experiment imposed no floor on
how far a maintainer may prune. An operator deploying the protocol should impose one.

\subsection{Lifting the comment strip}
\label{sec:appendix:failures}

Code separates the comment channel from the executor, and that design choice carries a measurable
price. We lift it in one run identical to the canonical one in every respect except that we render
the protocol arm's comments into the executor's prompt as well as the maintainer's, at one seed.
Exposure is real rather than nominal, since \InlineExposure{} of executor calls carry a prompt that
could show a comment at all, the rest being steps at which the prompt was still empty.

We pair per world against the canonical run over \InlineWorlds{} worlds: the prompt ends at
\InlineExcessPct{} excess size against \InlineRefExcessPct{} (\InlineDeltaExcessPp{},
\InlineDeltaExcessCI{}) at a satisfaction difference of \InlineDeltaSatPp{} (\InlineDeltaSatCI{},
straddling zero). Writing the
rationale into the file the model reads costs prompt size and buys no correctness. The effect sits
entirely in deletions rather than additions, and the direction reverses on a substantial minority
of worlds, so we report a direction at one seed rather than a settled magnitude. We have not tested
why an executor-visible comment should make a maintainer delete less.

Contamination of answer quality does not appear to be part of the cost. Maintenance vocabulary
appears in \InlineLeakage{} of the exposed arm's completions against \InlineLeakageControl{} under
the stripped prompt, so the executor does not visibly parrot the notes back. The strip earns its
place in this experiment by leaving the comment channel's information content as the only thing
that varies across arms, which licenses the causal reading in \autoref{sec:protocol}. We do not
claim it as a defense against leakage that did not occur.

